% =====================================================================
%  GMN2 §9.4  例子：N = 1, 2, 3
%  原文：wkb-paper.tex 6357–6845 行；印刷页 p.107–115；编号公式 (9.16)–(9.35) 共 20 条；无脚注
%  图：44 N2plot / 45 square-evolution / 46 N3plot / 47 penta-evolution / 48 pentaformula / 49 N3uplane
%  本块公式均在块内，故 \eqref 可用；跨块引用硬编码原文编号并加 % TODO
% =====================================================================

\subsection{例子}
\label{subsec:Examples}

现在我们用 $P(z)$ 为次数 $N=1,2,3,4$ 的多项式的那些情形，来演示本形式体系的各个侧面。一路上我们会遇到一些漂亮的「现实世界」例子，其中的穿墙公式只涉及有限多个 BPS 态。

由于我们在 $z=\infty$ 处有不规则奇点，要定义 WKB 三角剖分，就必须采用第 8 节% TODO 合稿改 \ref{sec:irregular}
的修改规则。把这些规则用到当前情形，我们得到的是曲面 $C'$ 的一个三角剖分；$C'$ 就是从 $\IC\IP^1$ 中挖去 $z=\infty$ 周围一个圆盘所得的曲面，其顶点是该圆盘上的标记点，对应于 WKB 曲线逃向 $z=\infty$ 的那些位置。正如我们在第 6 节% TODO 合稿改 \ref{sec:WKB-triangulation}
中一般性地解释过的，对某些特殊的 $\vartheta$ 值会出现有限 WKB 曲线，它们对应于我们场论的 BPS 态，并使 $T_\WKB(\vartheta)$ 发生跳变。


\subsubsection{$N=1$}

我们从 $N=1$ 开始，此时
\begin{equation}
P(z) = z. \label{eq:N1-P}
\end{equation}
% (9.16)
这里没有任何形变——无论可归一化与否——所以 $\CB$ 只是一个点。相应地，Hitchin 方程也只有唯一一个解，所以 $\CM$ 同样只是一个点。这个解可以显式写出：
\begin{equation}
\varphi = \begin{pmatrix} 0 &  \vert z \vert^{1/2} e^h \\
\frac{z}{\vert z \vert^{1/2}} e^{-h} & 0 \end{pmatrix}, \label{eq:ExactSol}
\end{equation}
% (9.17)
\begin{equation}
A = \left( \frac{1}{8} + \frac{1}{4} \vert z\vert \frac{d}{d\vert
z\vert} h \right) \sigma^3 \left( \frac{dz}{z}- \frac{d\bar z}{\bar
z} \right). \label{eq:ExactSol-A}
\end{equation}
% (9.18)
其中 $h(|z|)$ 是一个 Painlevé III 超越函数：令 $r := \frac{8 R}{3} \vert z\vert^{3/2}$，它满足
\begin{equation}
\left(\frac{d^2}{dr^2} + \frac{1}{r} \frac{d}{dr}\right)h = \half
\sinh(2h), \label{eq:Painleve}
\end{equation}
% (9.19)
边界条件为：当 $r\to 0$ 时 $h(r) \to - \frac{1}{3} \log r + const$。可以证明 \cite{McCoy:1976cd}，当 $r\to \infty$ 时 $h(r) \to \pi^{-1} K_0(r)$。下文第 13 节% TODO 合稿改 \ref{sec:R-To-Infinity}
中我们还会回到这个解。

$z=\infty$ 周围有 $N + 2 = 3$ 条 WKB 射线，由 (8.2) 给出，% TODO 跨块引用
而在 $z=0$ 处有一个转向点。$\CA$-平坦性方程多多少少有点像 Airy 方程：特别地，沿着三条射线 $r_k$ 各有一个小平坦截面 $s_k$。WKB 三角剖分 $T_\WKB(\vartheta)$ 只由一个三角形构成，它随 $\vartheta$ 的变化而在 $z$-平面中旋转。这里没有 flip，相应地也没有 BPS 态。这里没有 Fock–Goncharov 坐标，因为该三角剖分没有内部边。


\subsubsection{$N=2$}

接着考虑 $N=2$ 的情形。我们写
\begin{equation}
P(z) = z^2 + 2m, \label{eq:N2-P}
\end{equation}
% (9.20)
于是 $\Delta(z) = z + \frac{m}{z}$。参数 $m$ 是对数可归一化的。库仑分支 $\CB$ 是一个点，也没有 $U(1)$ 规范场。模空间 $\CM$ 同样是一个点。

谱曲线 $\Sigma$ 亏格为零，且在 $z = \infty$ 之上有两个穿刺点。此时 $H_1(\Sigma; \IZ)$ 是一维的，并且在两叶互换之下为奇，故 $\hat \Gamma = H_1(\Sigma; \IZ) \simeq \IZ$。这个一维格由单个味荷生成，$m$ 就是相应的质量参数。

在这个情形，我们不知道 Hitchin 方程的显式解。尽管如此，按照前几节的一般配方，让我们来考察 WKB 三角剖分。这里有 $N + 2 = 4$ 条 WKB 射线与 $N = 2$ 个转向点。WKB 三角剖分的一般行为如图 44 所示。从组合上看，四条边界边拼成一个正方形，而唯一一条内部边给出这个正方形的一个三角剖分。\insfig{3.3in}{N2plot}{$N=2$ 时一个一般的 WKB 三角剖分。分离 WKB 曲线以红色标出，而取作三角剖分各边的一般 WKB 曲线为黑色。阴影区域是我们在不规则奇点 $z = \infty$ 周围挖去的四个「花瓣」之并集。}

由于 $\hat\Gamma$ 是一维的，独立的达布坐标只有一个 $\CX_{\gamma}$。事实上我们断言它就等于 (9.13) 中给出的 $\mu$。% TODO 跨块引用
（确实，这是 $m, m^{(3)}$ 的、携带一个单位味荷的唯一合理函数。）为说明这一点，我们利用 $s_i$ 的渐近行为。由于两条 WKB 射线被一条斯托克斯射线隔开，在 $r_2$ 周围的一个角扇区内，$s_1$ 与 $s_3$ 都是指数式增长的，因而（若归一化得当）只相差小解 $s_2$ 的一个倍数，
\begin{equation}
s_1 = s_3 + a s_2 \label{eq:s1s3a}
\end{equation}
% (9.21)
其中 $a$ 是某个常数。类似地，在 $r_4$ 周围，$s_1$ 与 $s_3$ 都是指数式大的，因而（把形式单值、即渐近式的单值考虑进去之后）
\begin{equation}
s_1 = \mu s_3 + b s_4 \label{eq:s1mus3}
\end{equation}
% (9.22)
其中 $b$ 是某个常数。由于 $s_1, \dots, s_4$ 都是单值的，这些关系在整个平面上都成立，可用来把 Fock–Goncharov 坐标化简为
\begin{equation}
\CX_E =  -\frac{(s_1 \wedge s_2) (s_3 \wedge s_4)}{(s_2 \wedge s_3) (s_4
\wedge s_1)} = -\mu^{-1}. \label{eq:XE-N2}
\end{equation}
% (9.23)

当 $\vartheta$ 走过一段长度为 $\pi$ 的弧时，我们遇到一个临界值 $\vartheta = \vartheta_c$，在它那里出现一条连接两个转向点的有限 WKB 曲线。WKB 三角剖分 $T_\WKB(\vartheta)$ 在这个 $\vartheta$ 处经历一次 flip。见图 45。这次 flip 把唯一的 Fock–Goncharov 坐标变为 $\CX^T_{E} = 1/\CX^{T'}_{E}$（两式中 $E$ 都指那唯一一条内部边）。另一方面，$\CX^\vartheta_\gamma$ 保持不变，仍等于 $-\mu^{-1}$。

\insfig{5.5in}{square-evolution}{当相位 $\vartheta$ 从 $0$ 变到 $\pi$ 时，WKB 三角剖分的演化很简单，只发生一次 flip（虚线）。注意初、末两个三角剖分完全相同，只是 WKB 射线的标号有变。}

这一次 flip 意味着什么？它意味着，在上面讨论的标度极限下，我们保留了一个 BPS 粒子，其味荷为 $1$。基于这个局部模型，我们可以作一个更一般的预测：在库仑分支上最简单的那些奇点附近——即两个转向点相撞之处——会有一个 BPS 粒子变轻，它对应于张在两个转向点之间的一根 BPS 弦。


\subsubsection{插叙：$N=2$ 与周期 Taub-NUT 空间} \label{sec:pnut}

至此，我们可以与上一篇论文 \cite{Gaiotto:2008cd} 中考虑过的那个最根本的量子修正库仑分支的例子建立联系了。这将给出一个最近在胶子散射振幅的计算中用得上的结果 \cite{Alday:2009yn}。

在 $N=2$ 的例子中，我们只遇到一个对数可归一化的形变参数，相应地 $\CM$ 是平凡的——这一点起初也许会有点令人失望。人们原本也许会指望，能从中学到关于三维库仑分支上精确度量在四维理论奇点附近行为的某些信息。

从物理上，我们知道该期待什么。在原来的完整理论中，那个唯一的轻 BPS 粒子的味对称实际上是一个规范对称性。如果我们干脆忽略有质量 BPS 粒子的效应，那么 $\IR^3 \times S^1$ 上耦合单个轻超多重态的一个 $U(1)$ 规范理论的库仑分支，就是以「周期 Taub-NUT」度量著称的 \cite{Ooguri:1996me,Seiberg:1996ns}，它在 \cite{Gaiotto:2008cd} 的分析中起了核心作用。它并不是一个完备的超凯勒度量；毋宁说，它只定义到相差一个任意的质量标度 $\Lambda$。

在 \cite{Gaiotto:2008cd} 中，周期 Taub-NUT 度量是用一对全纯达布坐标 $(\CX_e, \CX_m)$ 来描述的。$\CX_e$ 与这里讨论的 $\CX_\gamma$ 一致。更有意思的坐标是 $\CX_m$，它经历斯托克斯现象，其斯托克斯因子 $\CK_\gamma$、$\CK_{-\gamma}$ 对应于单个有质量超多重态及其 CPT 共轭。

我们能否在当前的语境中为 $\CX_m$ 找到一个角色？在平坦联络模空间的层面上，很自然地去考虑让奇点处的单值参数 $\mu_a$ 变动。确实，$\mu_a$ 可以解释为第 $a$ 个奇点处剩余规范变换的复值矩映射。我们可以让 $\mu_a$ 变动，同时把规范变换限制为在奇点处趋于恒等。这就定义了一个更大的辛流形，每个奇点多出两个复坐标 \cite{MR1864833}。在我们当前的语境下，这意味着把 $\CM$ 从复维数 $0$ 扩大到复维数 $2$。正如我们马上会看到的，$\CX_e$ 与 $\CX_m$ 将实现为这个扩大后的 $\CM$ 上的函数。

一旦奇点处的规范变换被限制为恒等，为小的解选取一个特定的整体归一化就有规范不变的意义了，例如可以精确指定奇点处的渐近行为：
\begin{equation} \label{eq:gaugefix1}
\begin{split}
s_i(z,\zeta) \exp \Bigl[ &\frac{R}{\zeta}\Bigl(\frac{z^2}{2} + m \log z\Bigr)
-m^{(3)} (\log z - \log \bar z) \\
&+ R\zeta\Bigl(\frac{\bar z^2}{2} + \bar m \log \bar z\Bigr)\Bigr]
= \binom{1}{0} + O(1/z)
\end{split}
\end{equation}
% (9.24)
或
\begin{equation} \label{eq:gaugefix2}
\begin{split}
s_i(z,\zeta) \exp \Bigl[ - &\frac{R}{\zeta}\Bigl(\frac{z^2}{2} + m \log z\Bigr)
+m^{(3)} (\log z - \log \bar z) \\
&- R\zeta\Bigl(\frac{\bar z^2}{2} + \bar m \log \bar z\Bigr)\Bigr]
= \binom{0}{1} + O(1/z).
\end{split}
\end{equation}
% (9.25)

在我们现在考虑的这个例子里，这使我们能够对斯托克斯矩阵的单个元素——例如 $s_1 = s_3 + a s_2$ 中的 $a$——赋予规范不变的意义。（请注意，这个关系要与 \eqref{eq:gaugefix1}、\eqref{eq:gaugefix2} 相容，必须对 $\log z$ 的割线作恰当选取；我们把割线放在 $r_3$ 稍下方。）我们把 $a$ 看作一个「比值」（与「交比」相对），
\begin{equation}
a = \frac{s_1 \wedge s_3}{s_2 \wedge s_3}. \label{eq:ratio-a}
\end{equation}
% (9.26)

若 $\arg \zeta$ 的取值落在半平面 $\IH_\vartheta$ 内，而该 $\vartheta$ 使得 $r_1$ 与 $r_3$ 之间有一条边，我们就可以用 WKB 方法对小的 $\zeta$ 计算 $a$ 的渐近式。计算是直截了当的，给出我们惯常的指数形式 $a \sim c \exp (\pi R \zeta^{-1} Z_{\tilde\gamma})$，其中 $Z_{\tilde\gamma}$ 是 $\pi^{-1} \lambda$ 在圈 $\tilde \gamma$ 上的正则化周期，该圈从无穷远沿 $r_1$ 出发、沿 $r_2$ 回到无穷远，并绕过三角形 $123$ 中的那个转向点：
\begin{equation}
Z_{\tilde \gamma} = 2 \int_{z=\sqrt{-2m}}^L \sqrt{z^2 + 2 m}  - L^2-2m \log L = - m \log \frac{m}{-2e} + \CO\left(\frac{1}{L}\right). \label{eq:Zgamma-tilde}
\end{equation}
% (9.27)

这条路径与 $\gamma$ 相交一次。受这些渐近式的启发，我们对这个 $\vartheta$ 值定义 $\CX^\vartheta_{\tilde \gamma} = a$。

当把相位 $\vartheta$ 增大、越过使三角剖分翻转的临界值 $\vartheta_c$ 之后，我们应当转而考虑由 $s_4 = s_2 + a' s_1$ 定义的 $a'$，即
\begin{equation}
a' = \frac{s_4 \wedge s_2}{s_1 \wedge s_2}. \label{eq:aprime}
\end{equation}
% (9.28)

$a'$ 的 WKB 计算与 $a$ 的完全相同，只差一个至关重要的整体符号；所以这里的渐近式提示我们，应当用 $1/a'$ 来定义 $\CX^\vartheta_{\tilde \gamma}$。回忆关系式 $s_1 = s_3 + a s_2$ 与 $s_1 = \mu s_3 + b s_4$。按我们对 $\log z$ 所选的割线，若把这些关系一直延拓到 $r_1$，就会看到 $b s_4 = (1-\mu) s_1 + a \mu s_2$。与 $s_4 = s_2 + a' s_1$ 比较可知 $b = a \mu$，且 $a' = - a^{-1} (1-\mu^{-1})$。于是，若 ${\rm Re}(im/\zeta)>0$，则只要对这个 $\vartheta$ 定义 $\CX^\vartheta_{\tilde \gamma} = -\frac{1}{a'}$，$\CX_{\tilde \gamma}$ 的渐近行为就仍然自洽。

现在我们已经对翻转两侧的 $\vartheta$ 都定义了坐标 $\CX^\vartheta_\gamma$ 与 $\CX^\vartheta_{\tilde\gamma}$，而这次翻转与荷为 $\gamma$ 的 BPS 态相关联。由上面这些坐标的显式公式，很快可以看出：跨越这次 flip 的坐标变换恰好就是 $\CK_\gamma$！类似地，由荷为 $-\gamma$ 的 BPS 态所诱导的坐标变换结果是 $\CK_{-\gamma}$。在周期 Taub-NUT 的底空间上的坐标 $a^{PT}$ 与我们的质量参数之间取识别 $a^{PT} = -2 i m$，我们便看到 $\CX_\gamma, \CX_{\tilde \gamma}$ 具有与 $\CX_e, \CX_m$ 相同的渐近行为（在适当选取截断之后），并且在 $\vartheta$ 穿过 BPS 射线时以相同的方式变换。于是 \cite{Gaiotto:2008cd} 中黎曼–希尔伯特问题解的唯一性保证了 $\CX_\gamma, \CX_{\tilde \gamma}$ 与 $\CX_e, \CX_m$ 相重合。特别地，可以用 \cite{Gaiotto:2008cd} 的公式来计算 $\CX_{\tilde \gamma}$。

我们已经看到，与这个 Hitchin 系统相伴的辅助平坦性方程的斯托克斯数据，可以完全精确地用函数 $\CX_m(\zeta)$ 算出来——尽管我们算不出 Hitchin 方程的解！

到目前为止，我们通过把这个模空间看作平坦联络模空间的一个扩张（即引入一些与奇点相关的额外参数），说明了函数 $\CX_e$ 与 $\CX_m$ 在这个例子中如何出现。但我们还没有把这个模空间当作一个超凯勒流形来看待，因而尚未在 Hitchin 系统的语境中为周期 Taub-NUT 度量找到一个确切的位置。要做到这一点，我们可以试着让奇点处的三个质量参数 $m, m^{(3)}, \bar{m}$ 变动，然后再只考虑在奇点处化归为恒等的规范变换，从而加上一个取值为圆的参数。然而，我们并不清楚能否在所得到的扩张模空间上定义一个超凯勒度量；度量和超凯勒形式在作用到质量的变分上时会发散。人们可以试着挖去每个奇点周围一个小圆盘来正则化这个发散。度量将以对数方式依赖于截断半径，与周期 Taub-NUT 度量的情形一模一样，并且可能是不完备的。由于我们对超凯勒度量的扭量构造在底空间上是完全局部的，这种可能的不完备性并不构成障碍：度量可以按我们的标准做法由坐标 $\CX_\gamma, \CX_{\tilde \gamma}$ 算出；而由于它们与 $\CX_e, \CX_m$ 重合，该度量就将与周期 Taub-NUT 度量重合！

这条进路可以推广到更复杂的 Hitchin 系统，但本文不再进一步追究。


\subsubsection{$N=3$} \label{sec:N3}

现在我们来到 $N=3$ 的情形，这里我们将第一次遇到一个穿墙公式。我们写
\begin{equation} \label{eq:N3}
P(z) = z^3 - 3 \Lambda^2 z + u,
\end{equation}
% (9.29)
其中 $\Lambda$ 是定义该理论的一个不可归一化参数，$u$ 是参数化库仑分支 $\CB$ 的一个可归一化模。$P(z)$ 的判别式为 $27((2\Lambda^3)^2-u^2)$，所以 $\CB$ 上有两个奇点，位于 $u = \pm 2 \Lambda^3$，在那里 $P(z)$ 的两个零点相撞。（我们取 $\Lambda \neq 0$，这样就不存在使三个零点全部相撞的 $u$。）

照例，$\CB$ 上有一个格的局部系统 $\hat\Gamma$，它由带穿刺点的椭圆曲线族 $\Sigma_u$ 的同调之奇部给出。此时 $\hat\Gamma$ 的秩为 $2$。总而言之，$\CB$ 与 $SU(2)$ Seiberg–Witten 理论在 $N_f = 0$ 时那个众所周知的 $u$-平面非常相似（下文第 10.1 节% TODO 合稿改 \ref{sec:nf0}
中会考虑它）。

\insfig{2.7in}{N3plot}{$N = 3$ 例子中一个典型的 WKB 三角剖分，取一般的 $\vartheta$。}

现在让我们考察 WKB 三角剖分。$z = \infty$ 处有五条 WKB 射线（相应地有五条斯托克斯射线），另有三个转向点。一般的 $T_\WKB(\vartheta, u)$ 是一个五边形的三角剖分，如图 46 所示。取与这个三角剖分的两条内部边 $E_{1,2}$ 相关联的 $\gamma_E^\vartheta$，我们就得到 $\hat\Gamma_u$ 的一组整基。具体说来，让我们把 $\{\gamma_1, \gamma_2\}$ 定义为由 $T_\WKB(\vartheta = 0, u = 0)$ 如此得到的基，并取 $\langle \gamma_1, \gamma_2 \rangle = 1$。

当固定 $u$、让 $\vartheta$ 从 $0$ 变到 $2 \pi$ 时，$T_\WKB(\vartheta,u)$ 在周期 $\pm Z_{\gamma_1}$ 与 $\pm Z_{\gamma_2}$ 的相位处共经历 $4$ 次 flip 跳变。如果只让 $\vartheta$ 走过一段长为 $\pi$ 的范围，那么 $T_\WKB(\vartheta,u)$ 翻转两次。这两次 flip 由两条独立的有限 WKB 曲线生成，代表两个 BPS 粒子（另外两次 flip 对应于它们的反粒子）。见图 47 的上栏。\insfig{5.5in}{penta-evolution}{上栏：$u=0$ 时，随相位 $\vartheta$ 从 $0$ 变到 $\pi$，$T_\WKB(\vartheta, u)$ 的演化很简单，只发生两次 flip（虚线）。下栏：当 $u$ 越过边缘稳定墙之外时，$T_\WKB$ 的演化方式不同，涉及三次 flip（虚线）。注意 $T_\WKB(\vartheta) = T_\WKB(\vartheta + \pi)$，只差 WKB 射线的重新标号，而后者等价于一次 pop 变换。} 在 $\CB$ 的每个奇点处，这两个 BPS 态之一会变成无质量的。（正如我们上面提到的，这些奇点出现在一对转向点相撞之时；那个无质量的 BPS 态于是对应于连接这对转向点的一条有限 WKB 曲线。）

在 $\CB$ 中有一条闭合的边缘稳定墙，它把 $u$-平面分成两个连通分支（同样，很像 $SU(2)$ Seiberg–Witten 理论中的情形）。到目前为止我们讨论的是 $u=0$，它位于墙内。若取某个位于墙外的 $u$，$T_\WKB(\vartheta, u)$ 的行为就不同：此时有三次 flip，由三条有限 WKB 曲线诱导，对应于三个 BPS 态。见图 47 的下栏。（取 $u$ 很大时最容易看出这一点。）

更精确地说，两个奇点 $u = \pm 2 \Lambda^3$ 把边缘稳定墙分成两段。沿其中一段，$Z_{\gamma_{1,2}}$ 的相位对齐。从墙内穿到墙外、穿过这一段之后，我们发现三个 BPS 态，其荷为 $\gamma_1, \gamma_2, \gamma_1 + \gamma_2$。若穿过墙的另一段——那里 $Z_{\gamma_1}$ 与 $Z_{-\gamma_2}$ 对齐——我们得到的则是荷为 $\gamma_1, \gamma_2, \gamma_1-\gamma_2$ 的 BPS 态。为看出这两个结果彼此相容，我们回忆：$\hat\Gamma$ 绕两个点 $u = \pm 2 \Lambda^3$（$P(z)$ 的两个零点在那里相撞）具有标准的 Lefschetz 单值。绕这两个奇点的顺时针单值变换的生成元可以取为：
\begin{equation}\label{eq:LocalSys-1}
\begin{aligned}
M_1 &= (\gamma_1, \gamma_2) \to (\gamma_1, \gamma_2-\gamma_1), \\
M_2 &= (\gamma_1, \gamma_2) \to (\gamma_1+\gamma_2, \gamma_2).
\end{aligned}
\end{equation}
% (9.30)

这三条有限 WKB 曲线在我们走向任意大的 $|u|$ 时都保持在谱中（确实，没有别的边缘稳定墙能让它们消失）。由此我们推断：在更一般的理论中，在库仑分支上三个转向点合并的那种奇点附近，通常会存在三个轻的 BPS 粒子，它们实现为连接这三个转向点的 BPS 弦。

\insfig{4.5in}{pentaformula}{在墙的两侧，从 $\vartheta=0$ 的三角剖分走到 $\vartheta=\pi$ 的三角剖分有两条不同的路径，它们必定对应同一个辛同胚。五边形恒等式由此得出。}

\subsubsection*{穿墙公式}

现在可以看穿墙公式如何浮现了。设 $u_-$ 是墙内一点、$u_+$ 是墙外一点，并固定某个 $\vartheta$。我们在 $\CM$ 上有两个坐标系 $\CX^{\vartheta, u_-}_\gamma$ 与 $\CX^{\vartheta+\pi, u_+}_\gamma$，它们由某个辛同胚相联系。为计算这个辛同胚，我们考察三角剖分 $T_\WKB(\vartheta, u)$ 沿从 $(\vartheta, u_-)$ 到 $(\vartheta+\pi, u_+)$ 的一条路径的演化。

一种可能是：先把 $\vartheta$ 变到 $\vartheta+\pi$，再从 $u_-$ 形变到 $u_+$。当 $\vartheta$ 变化时，三角剖分经历两次 flip，对应于墙内的两个 BPS 态。当我们变动 $u$ 时，只要没有 $\arg Z_{\gamma}(u)$ 穿过 $\vartheta$（这一点总可以通过取 $u_\pm$ 离墙足够近、$\vartheta$ 足够一般来做到），三角剖分根本不跳变。所以三角剖分的总变换恰好涉及两次 flip。

然而还有另一种可能：先从 $u_-$ 形变到 $u_+$，再把 $\vartheta$ 变到 $\vartheta+\pi$。那样的话，三角剖分经历三次 flip，对应于墙外的三个 BPS 态。

这两种算法必定给出同一个辛同胚。这对应于恒等式
\begin{equation}\label{eq:simple-pentagon}
\CK_{\gamma_1} \CK_{\gamma_2} = \CK_{\gamma_2} \CK_{\gamma_1+\gamma_2} \CK_{\gamma_1}
\end{equation}
% (9.31)
它对墙的其中一段成立；而在另一段上则是
\begin{equation}\label{eq:simple-pentagon-2}
\CK_{-\gamma_2} \CK_{\gamma_1} = \CK_{\gamma_1} \CK_{\gamma_1-\gamma_2} \CK_{-\gamma_2}
\end{equation}
% (9.32)
人们可以直接计算验证，这些关系确实成立。这两条基本恒等式中的任何一条，我们都将称为「五边形恒等式」。

\insfig{4.0in}{N3uplane}{$N=3$ 理论的库仑分支 $\CB$，其中标出了 $Z_{\pm \gamma_1}$ 与 $Z_{\gamma_2}$ 对齐的边缘稳定墙。注意 $\hat\Gamma$ 有单值，因此为了整体地定义基 $\{\gamma_1, \gamma_2\}$ 我们必须引入割线，图中以橙色标出。在每个区域内我们给出 BPS 谱以及 BPS 射线的循环次序。沿着墙则给出穿墙公式。}

我们在 $\CB$ 上讨论过的所有结构都图示在图 49 中。


\subsubsection{插叙：辛同胚与单值}

我们上面遇到的这些辛同胚，与绕 $\CB$ 中一个奇点转一圈时出现的单值变换之间，有一个有趣的关系。

考虑当 $u$ 绕过最简单的一类奇点时 BPS 射线的行为；在这类奇点处，某个 $Z_{\gamma_0}(u)$ 有单零点。这正是我们上面在 $N=2$ 例子中凑近观察过的那类奇点。当 $u$ 绕该奇点一圈时，$Z_{\gamma_0}(u)$ 的相位转过 $2 \pi$。于是 BPS 射线 $\ell_{\gamma_0}$ 会从所有其他 BPS 射线上扫过，随后是射线 $-\ell_{\gamma_0}$ 扫过。换言之，绕该奇点一圈，我们就穿过了一连串的边缘稳定墙。

容易验证 $\CK_{\gamma_0} \CK_{-\gamma_0}$ 就是变换
\begin{equation}
   \CX_\gamma \to (-\sigma(\gamma_0))^{\langle \gamma,\gamma_0 \rangle}
   \CX_{\gamma+ \langle \gamma,\gamma_0 \rangle \gamma_0}. \label{eq:KS-monodromy}
\end{equation}
% (9.33)
于是，当 BPS 射线 $\ell_{\gamma_0}$ 与 $\ell_{-\gamma_0}$ 扫过整个谱时，相应的 Kontsevich–Soibelman 变换实现的正是所预期的荷的单值变换。（注意，二次加细 $\sigma$ 可以取得使：只要 $\gamma_0$ 是某个 BPS 超多重态的荷，就有 $\sigma(\gamma_0) = -1$，这一点很方便；这样的二次加细在这个单值变换下不变，因而后者只作用在 $\CX_\gamma$ 的荷标号上。）

那么三个转向点合并的那种奇点又如何呢，就像我们在 $N=3$ 情形中研究过的？如上面所指出，在这种奇点附近通常会出现三个轻的 BPS 粒子，实现为连接各转向点的 BPS 弦。它们的荷投影到与标度极限相关的规范荷二维格上，形如 $\gamma_1$、$\gamma_1 + \gamma_2$、$\gamma_2$。天真地看，绕奇点转圈时人们也许会预期，每条 BPS 线都会被三个轻粒子的 BPS 线扫过，然后再被三个轻反粒子的 BPS 线扫过。但这其实是不对的，只要对 KS 变换作一点简单操作就能明白。确实，我们有
\begin{equation}\label{eq:KS-chain}
\begin{split}
&\CK_{\gamma_2}\CK_{\gamma_1+\gamma_2}\CK_{\gamma_1} \CK_{-\gamma_2}\CK_{-\gamma_1-\gamma_2}\CK_{-\gamma_1}
= \CK_{\gamma_1}\CK_{\gamma_2} \CK_{-\gamma_1}\CK_{-\gamma_2} \\
&\quad = (\CK_{\gamma_1} \CK_{-\gamma_1})\CK_{-\gamma_1+\gamma_2}(\CK_{\gamma_2}\CK_{-\gamma_2})
= (\CK_{\gamma_1} \CK_{-\gamma_1})(\CK_{\gamma_2}\CK_{-\gamma_2})\CK_{-\gamma_1}.
\end{split}
\end{equation}
% (9.34)

这并不完全是所预期的单值作用——后者本来只由 $(\CK_{\gamma_1} \CK_{-\gamma_1})(\CK_{\gamma_2}\CK_{-\gamma_2})$ 实现。为化解这个困难，应当更仔细地考察那些小的中心荷对 $u$ 的精确依赖。为简单起见取 $\Lambda \to 0$。于是由 \eqref{eq:N3} 可以看出，六个轻 BPS 粒子的中心荷行为如同六个根 $u^{5/6}$。由此可见，绕 $u=0$ 一圈只把中心荷旋转 $\exp \frac{5\pi i}{3}$。因此绕这一圈时，一条一般的 BPS 射线经历的是五次穿墙而非六次，而相应的恒等式实际上是
\begin{equation}
\CK_{\gamma_2}\CK_{\gamma_1+\gamma_2}\CK_{\gamma_1} \CK_{-\gamma_2}\CK_{-\gamma_1-\gamma_2}
= (\CK_{\gamma_1} \CK_{-\gamma_1})(\CK_{\gamma_2}\CK_{-\gamma_2}), \label{eq:KS-chain-final}
\end{equation}
% (9.35)
它确实给出所要的单值变换。
